\section {Metric of Success for CFP and Backup Plans} \label{sec:backup}
As a baseline metric of success, we expect at least 10-15 paper submissions based on $12$ total submissions in the last year's workshop.  We anticipate 6-8 accepted papers at the workshop that relate to the workshop goals (Sec.~\ref{sec:workshopGoals}). Another metric of success for CFP would be submissions from scientists with diverse backgrounds (mathematics, application science, psychology) in addition to core visualization researchers on identified challenges in uncertainty visualization.

In the case where we only get a few submissions or do not find expected number of panelists, we will compensate these sessions by inviting renowned keynote speakers who can discuss challenges and potential solutions relevant to each of the mentioned themes ($\mathbf{T_1}$ and $\mathbf{T_2}$). We will again stress more on Q\&A session at the end of keynote talk to foster interaction and meet workshop goals. 


%Another potential idea is to arrange breakout sessions based on the main challenges identified in the proposed panel sessions {\em"Top Challenges in Uncertainty Visualization"}. We will then write a position paper through ideas discussed in the breakout sessions. For breakout sessions however we will need to have access to separate rooms or arrange a way to separate groups to avoid any audibility conflicts with other groups. 


%As a baseline metric of success, we expect at least 8 submissions for papers and 8 submissions for posters. Another metric of a success for CFP would be submissions from scientists with diverse backgrounds (mathematics, application science, psychology) in addition to core visualization researchers on preliminary advances, challenges, or use cases in uncertainty analysis.

%In the case, we only get a few number of accepted submissions or no submissions for papers/posters, we will compensate the presentation time with external invited speakers or another panel session. In the case of another panel session, we plan to discuss the topic of {\em "Applications and Uncertainty"}, in which we will invite domain-experts from various application areas who can discuss ways to understand and deal with uncertainty in their domain. We will invite one/two external speakers as the other backup idea who can talk about the importance and treatment of uncertainty in their domain.